Only one assignment is observed from each schedule.  Consequently \(P_0\) and \(\bar P_0=\bigotimes_zP_{0,z}\) have the same one-cluster observed law, which proves \(Q^B_{\bar P_0,p}=Q^B_{P_0,p}\).  Assumptions \(\text{ass:iid-schedules}\), \(\text{ass:bernoulli-label-iid}\), \(\text{ass:bernoulli-label-schedule-independence}\), and \(\text{ass:bernoulli-units}\) make those records independent and give the exact-hit probabilities in \(\text{lem:exact-hit-factorization}\).

Fix any \(h\in\mathbb R^A\); no centering restriction is imposed.  On \(z\in\mathcal S_k\), differentiation of the finite-support exponential family gives score
\[
h_k\{W(z)-\mu_z\}/\tau_k^2
\]
and
\[
\left.\frac d{dt}\mathbb E_{P_{t,h,z}}W(z)\right|_{t=0}
=\frac{h_k}{\tau_k^2}\operatorname{Var}_{P_0}\{W(z)\}.
\]
The division is valid because \(\tau_k^2\geq\underline\tau^2>0\) on \(A\).  Averaging uniformly over \(\mathcal S_k\) and using the definition of \(\tau_k^2\) yields
\[
\left.\frac d{dt}U_k(\bar P_{t,h})\right|_{t=0}
=\frac{h_k}{M_k\tau_k^2}\sum_{z\in\mathcal S_k}\operatorname{Var}_{P_0}\{W(z)\}=h_k. \tag{5}
\]
The observed score coordinate is therefore \(\dot\ell_k^B\).  It is centered, bounded, and, because distinct count slices are disjoint,
\[
\mathbb E_0(\dot\ell_j^B\dot\ell_k^B)
=\mathbf 1\{j=k\}\frac{q_k}{M_k\tau_k^4}
\sum_{z\in\mathcal S_k}\operatorname{Var}_{P_0}\{W(z)\}
=\mathbf 1\{j=k\}\frac{q_k}{\tau_k^2}. \tag{6}
\]
This is \(I_B\), and it includes the common-shift score \(\sum_k\dot\ell_k^B\) rather than deleting it.

The four-node early check is exact.  On the count-one slice, let all four unit outcomes copy one Bernoulli bit.  If every assignment-specific bit has mean \(1/2\), then \(\tau_1^2=1/4\), so (6) gives information \(4q_1\) and influence variance \(1/(4q_1)\).  If the four means are \((4/5,4/5,1/5,1/5)\), every cell variance is \(4/25\), hence \(\tau_1^2=4/25\), information \(25q_1/4\), and influence variance \(4/(25q_1)\).  Thus assignment heterogeneity changes the calibrated variance primitive but does not change the exact-hit multiplier.  At finite \(C\), the prescribed empty-cell fallback contributes bias
\[
\frac1{M_k}\sum_{z\in\mathcal S_k}(1/2-\mu_z)
\left(1-\frac{q_k}{M_k}\right)^C,
\]
which need not vanish identically but decays exponentially because \(q_k/M_k>0\); no finite-sample unbiasedness is used below.

Since \(0\leq W\leq1\), Taylor's theorem on the finite support gives, uniformly for \(a\) in a fixed neighborhood of zero,
\[
\Lambda_{z,k}(a)=\frac{a^2}{2}\operatorname{Var}_{P_0}\{W(z)\}+O(|a|^3). \tag{7}
\]
Substitute \(a=C^{-1/2}h_k/\tau_k^2\), sum over clusters, and apply the bounded law of large numbers to the assignment-cell frequencies.  The linear term is \(h^{\prime}\Delta_C^B\), (6)--(7) give the quadratic term \(-h^{\prime}I_Bh/2+o_P(1)\), and the total cubic remainder is \(O(C^{-1/2}\max_k|h_k|^3)\) for every fixed bounded set of representatives.  The summands in \(\Delta_C^B\) are iid centered bounded vectors with covariance \(I_B\).  Expanding their characteristic functions at the origin proves \(\Delta_C^B\Rightarrow N(0,I_B)\), and hence the stated LAN expansion.

Equation (5), the bounded third derivative, and \(\text{thm:efficient-exact-hit-bernoulli}\) give under every fixed local alternative, jointly on finite sets and uniformly on compact sets of representatives,
\[
\sqrt C\{\widetilde U_{C,A}-U_A(P_0)\}\Rightarrow
h+N(0,\Sigma_{B,A}). \tag{8}
\]
For the decision lower bound, fix finitely many representatives whose quotient norms are at most \(M\), allowing arbitrary finite common shifts, and put any finite prior on them.  Joint convergence of their likelihood ratios follows from LAN; bounded exponential-family scores give uniform integrability.  Therefore the limiting Bayes risk of any rule sequence is at least the Bayes risk in the full Gaussian shift experiment \(X\sim N(h,\Sigma_{B,A})\).  The scaled causal loss converges on that finite set to \(\max_kh_k-h_a\) by (5).  A fixed baseline gap makes any inactive action no better and, if chosen with nonvanishing probability, asymptotically worse.

Taking the supremum over finite priors gives the compact quotient minimax value.  To see that the common nuisance is included in this passage, start with a finite quotient prior and tensor it with a uniform grid of common shifts in \([-T,T]\).  Gaussian total-variation continuity lets the translation mesh tend to zero, and the boundary-to-volume ratio of this grid tends to zero as \(T\to\infty\).  The resulting Bayes lower values converge to the invariant-rule value produced by the same averaging calculation as in \(\text{lem:gaussian-quotient-reduction}\).  On the compact quotient ball, separating the convex set of finite risk vectors and using their uniform Gaussian total-variation continuity shows that finite quotient priors are sufficient.  Thus the common shift has not been declared known, invariant rules depend only on \(D_AX\), and the lower bound is \(G_{A,M}(\tau_A^2,q_A)\).  Finally that lemma gives \(G_{A,M}\uparrow G_A\), and \(\text{thm:nuisance-deletion-bernoulli}\) transfers the lower bound from the hit experiment to rules observing the full record.